\section{Acoustics}
\label{sec:acoustics}

For a machine that hovers beside a person's head, the acoustic model is not a
secondary consideration; \cref{sec:results} finds it to be the binding
constraint. It is also the place where the most error was found, so it is
developed carefully.

\subsection{Why the sixth power}

\begin{model}[Broadband loading noise]
\label{mod:acoustic}
At the scales and Reynolds numbers considered, rotor noise is dominated by
broadband loading noise: unsteady lift on the blade sections, arising from
inflow turbulence and boundary-layer vortex shedding, radiating as a
distribution of acoustic dipoles.
\end{model}

Curle's extension of Lighthill's acoustic analogy to flows bounded by solid
surfaces~\cite{lighthill1952,curle1955} shows that the surface term is dipole
in character, and dimensional analysis of the dipole source strength gives the
acoustic power radiated by a body of characteristic area $A_b$ in a flow of
characteristic velocity $U$:
\begin{equation}
  W_{\mathrm{ac}} \;\sim\; \frac{\rho\, A_b\, U^{6}}{c_0^{3}} ,
  \label{eq:dipole}
\end{equation}
with $c_0$ the speed of sound. This is the classical $U^6$ dipole law, against
$U^8$ for the quadrupole (free jet) case; see Ffowcs Williams and
Hawkings~\cite{fwh1969} for the rotating-surface generalisation, and
Brooks, Pope and Marcolini~\cite{brooks1989} for the airfoil self-noise
mechanisms this lumps together.

Taking $10\log_{10}$ of \cref{eq:dipole} with $U=\Vtip$ gives a
$60\log_{10}\Vtip$ dependence. Loading enters through the sectional lift, for
which we take an empirical $10\log_{10}T$, and $N$ statistically independent
rotors add incoherently, contributing $10\log_{10}N$. Hence

\begin{model}[Source model]
\label{mod:spl}
\begin{equation}
  L_p(\SI{1}{\metre})
  = C + 60\log_{10}\!\frac{\Vtip}{\Vtip^{\mathrm{ref}}}
      + 10\log_{10}\!\frac{T_{\mathrm{rotor}}}{T^{\mathrm{ref}}}
      + 10\log_{10} N ,
  \label{eq:splmodel}
\end{equation}
with $\Vtip^{\mathrm{ref}}=\SI{100}{\metre\per\second}$,
$T^{\mathrm{ref}}=\SI{5}{\newton}$ and $C$ an anchor calibrated against
measured aircraft in \cref{sec:verification}.
\end{model}

\begin{observation}[The design lever]
\label{obs:sixthpower}
Halving tip speed reduces the level by $60\log_{10}2 = \SI{18.1}{\decibel}$.
No other parameter in the model has comparable authority. Together with
\cref{cor:designrule}, which says that raising solidity lowers the tip speed
needed to carry a given thrust, this determines the entire shape of the
machine: large, fat-bladed, slowly turning rotors.
\end{observation}

\begin{figure}[tbp]
\centering
% Generated by paper/figures/gen_c.py. Do not edit by hand.
\def\cM{1.724}
\def\cMbat{0.288}
\def\cL{0.3827}
\def\cR{0.2460}
\def\cReach{0.629}
\def\cSpan{1.257}
\def\cKgap{1.100}
\def\cGap{0.0492}
\def\cSoverD{0.100}
\def\cKint{1.0252}
\def\cRcp{0.130}
\def\cCoMz{0.0275}
\def\cScreenOffz{0.1025}
\def\cScreenW{0.345}
\def\cScreenH{0.194}
\def\cArmr{8.0}
\def\cWall{0.66}
\def\cWallStress{0.083}
\def\cWallStiff{0.660}
\def\cWallMin{0.60}
\def\cDriver{stiffness}
\def\cSigmaWork{44}
\def\cSigmaAllow{350}
\def\cFbend{33}
\def\cFlimit{15.2}
\def\cZrotor{1.37}
\def\cZceil{1.03}
\def\cZoverR{5.57}
\def\cZcOverR{4.19}
\def\cKground{1.0020}
\def\cKceiling{1.0072}
\def\cKc{2.0}
\def\cCeilPct{0.7}
\def\cGroundPct{0.20}
\def\cKcSpanG{13}
\def\cCeiling{2.4}
\def\cBpf{34.5}
\def\cRpm{1036}
\def\cNb{2}
\def\cAbpf{-37.5}
\def\cLwa{61.8}
\def\cRroom{20.8}
\def\cRc{0.91}
\def\cStandoff{1.20}
\def\cSplEar{56.6}
\def\cSplOne{53.8}
\def\cSplFar{54.8}
\def\cTrotor{4.23}
\def\cVtip{26.7}
\def\cAnchor{82.9}
\def\cVi{3.01}
\def\cTfCross{21}
\def\cTf{30}
\def\cQmax{0.289}
\def\cMmot{61}
\def\cKtau{5.5}
\def\cEmot{0.30}
\def\cRoomL{5}
\def\cRoomW{4}
\def\cAlphaRoom{0.20}
%
% The source model at the design's thrust per rotor: 60 dB per decade of tip
% speed, with the design and a camera-drone tip speed marked.
\begin{tikzpicture}
  \begin{semilogxaxis}[paperplot, height=5.6cm,
      xlabel={tip speed $\Vtip$ [\si{\metre\per\second}]},
      ylabel={$L_p(\SI{1}{\metre})$ [\si{\decibel A}]},
      xmin=10, xmax=150, ymin=25, ymax=95,
      xtick={10,20,50,100,150}, xticklabels={10,20,50,100,150},
      legend pos=north west]
    \addplot[cA, domain=10:150, samples=60]
      {\cAnchor + 60*log10(x/100) + 10*log10(\cTrotor/5) + 10*log10(4)};
    \addlegendentry{\cref{eq:splmodel}, $T_{\mathrm{rotor}}=\SI{\cTrotor}{\newton}$, $N=4$}
    % comfort bands
    \draw[cC, dashed, line width=0.6pt] (axis cs:10,45) -- (axis cs:150,45);
    \node[lbl, cC, anchor=south east] at (axis cs:150,45) {quiet office};
    \draw[cD, dashed, line width=0.6pt] (axis cs:10,60) -- (axis cs:150,60);
    \node[lbl, cD, anchor=south east] at (axis cs:150,60) {speech};
    % halving the tip speed
    \pgfmathsetmacro{\Ldes}{\cAnchor + 60*log10(\cVtip/100) + 10*log10(\cTrotor/5) + 10*log10(4)}
    \pgfmathsetmacro{\Vtwo}{2*\cVtip}
    \pgfmathsetmacro{\Ltwo}{\Ldes + 60*log10(2)}
    \draw[cG, line width=0.5pt] (axis cs:\cVtip,\Ldes) -- (axis cs:\Vtwo,\Ldes) -- (axis cs:\Vtwo,\Ltwo);
    \node[lbl, cG, anchor=west] at (axis cs:\Vtwo,0.5*\Ldes+0.5*\Ltwo) {\SI{18.1}{\decibel}};
    \node[lbl, cG, anchor=north] at (axis cs:1.41*\cVtip,\Ldes) {$\times2$};
    \addplot[only marks, mark=*, mark size=2.3pt, cA] coordinates {(\cVtip,\Ldes)};
    \node[lbl, cA, anchor=north west] at (axis cs:\cVtip,\Ldes-1) {design};
    \pgfmathsetmacro{\Ldrone}{\cAnchor + 60*log10(68/100) + 10*log10(\cTrotor/5) + 10*log10(4)}
    \addplot[only marks, mark=square*, mark size=2.2pt, cB] coordinates {(68,\Ldrone)};
    \node[lbl, cB, anchor=south east, align=right] at (axis cs:66,\Ldrone+1)
      {camera drone,\\same thrust};
  \end{semilogxaxis}
\end{tikzpicture}

\caption{The source model \cref{eq:splmodel} at the design's thrust per rotor.
On a logarithmic tip-speed axis the level is a straight line of
\SI{60}{\decibel} per decade, so halving the tip speed buys \SI{18.1}{\decibel}
wherever one starts. The design turns at \SI{\cVtip}{\metre\per\second};
carrying the same thrust at a camera drone's tip speed would cost about
\SI{24}{\decibel}, the difference between a quiet office and a raised voice.}
\label{fig:sixth}
\end{figure}

\Cref{fig:sixth} draws \cref{eq:splmodel} at the design's thrust.

\subsection{Sound power against sound pressure}

\begin{definition}
The sound power level $L_W$ of a source is
$L_W = 10\log_{10}(W/W_0)$, $W_0=\SI{1e-12}{\watt}$, a property of the source
alone. The sound pressure level $L_p$ at a point is
$L_p = 20\log_{10}(p_{\mathrm{rms}}/p_0)$, $p_0=\SI{20}{\micro\pascal}$, and
depends on the environment.
\end{definition}

\begin{proposition}[Free-field conversion]
\label{prop:lwlp}
For a point source of directivity factor $Q$ radiating into free space above a
rigid plane, the mean-square pressure at distance $r$ satisfies
\begin{equation}
  L_p(r) = L_W + 10\log_{10}\!\left(\frac{Q}{4\pi r^2}\right) ,
  \label{eq:lwlp}
\end{equation}
so for a source above a reflecting plane, $Q=2$, and at $r=\SI{1}{\metre}$,
\begin{equation}
  L_p(\SI{1}{\metre}) = L_W - 10\log_{10}(2\pi) = L_W - \SI{7.98}{\decibel} .
  \label{eq:eightdb}
\end{equation}
\end{proposition}

\begin{proof}
Acoustic intensity at distance $r$ for a source of power $W$ radiating with
directivity $Q$ is $I = QW/(4\pi r^2)$. In the far field of a plane wave
$I = p_{\mathrm{rms}}^2/(\rho c_0)$, and the reference quantities are chosen so
that $10\log_{10}(I/I_0) = L_p$ with $I_0 = p_0^2/(\rho c_0)$ to within
\SI{0.1}{\decibel} at standard conditions. Taking logarithms gives
\cref{eq:lwlp}; substituting $Q=2$, $r=1$ gives \cref{eq:eightdb}.
\end{proof}

\Cref{fig:hemisphere} is the geometry behind $Q=2$.

\begin{figure}[tbp]
\centering
% A source on a reflecting floor radiates into a hemisphere of area 2 pi r^2,
% which is the Q = 2 of eq:lwlp and the source of the 8 dB between a declared
% sound power and a pressure at one metre.
\begin{tikzpicture}[>={Stealth[length=1.8mm]}, font=\footnotesize]
  \fill[pattern=north east lines, pattern color=cG] (-3.2,-0.25) rectangle (3.2,0);
  \draw[line width=0.8pt] (-3.2,0) -- (3.2,0);
  \node[anchor=north east, cG] at (3.2,-0.28) {reflecting plane};
  \foreach \r/\op in {1.0/70, 1.8/50, 2.6/32}
    \draw[cA!\op!white, line width=0.9pt] (\r,0) arc (0:180:\r);
  \foreach \ang in {20,55,90,125,160}
    \draw[vec, cA] (\ang:0.22) -- (\ang:0.82);
  \fill[cD] (0,0) circle (0.07);
  \node[anchor=north east, cD] at (-0.1,-0.28) {source of power $W$};
  \draw[dim, black] (0,0) -- (35:2.6) node[pos=0.66, above left=-2pt] {$r$};
  \node[anchor=south, cA] at (0,2.62) {wavefront, area $2\pi r^{2}$};
  \node[anchor=west, align=left] at (3.2,1.55)
    {$I=\dfrac{W}{2\pi r^{2}}=\dfrac{QW}{4\pi r^{2}},\quad Q=2$\\[6pt]
     $L_p(r)=L_W+10\log_{10}\dfrac{Q}{4\pi r^{2}}$\\[6pt]
     $L_p(\SI{1}{\metre})=L_W-10\log_{10}2\pi$\\[2pt]
     $\phantom{L_p(\SI{1}{\metre})}=L_W-\SI{7.98}{\decibel}$};
\end{tikzpicture}

\caption{Sound power against sound pressure, \cref{prop:lwlp}. A source on a
reflecting plane spreads its power $W$ over hemispherical wavefronts of area
$2\pi r^2$, which is the directivity $Q=2$ of \cref{eq:lwlp}. At one metre that
area is $2\pi\,\si{\metre\squared}$, and $10\log_{10}2\pi=\SI{7.98}{\decibel}$
is the whole difference between a declared sound power level and the pressure
level a listener at one metre hears.}
\label{fig:hemisphere}
\end{figure}

\begin{remark}[A trap worth naming]
European regulation 2019/945 requires manufacturers of unmanned aircraft to
declare a \emph{sound power} level $L_{WA}$. It is routinely misread as a
pressure at one metre. By \cref{prop:lwlp} that misreading is worth
\SI{8}{\decibel}, and in the wrong direction: it makes published aircraft sound
louder than they are, and so makes a design calibrated against the misreading
look quieter than it is. \Cref{sec:verification} records the consequences.
\end{remark}

\subsection{The room}

Free field is the wrong model indoors. The classical treatment
(Sabine; see Beranek~\cite{beranek1988} or Kuttruff~\cite{kuttruff2009}) writes
the steady-state level as the sum of a direct field, which falls with distance,
and a diffuse reverberant field, which does not.

\begin{definition}[Room constant]
For a room of total interior surface area $S$ and average absorption
coefficient $\bar\alpha$,
\begin{equation}
  \mathcal{R}_{\mathrm{room}} \coloneqq \frac{S\,\bar\alpha}{1-\bar\alpha} .
  \label{eq:roomconstant}
\end{equation}
\end{definition}

\begin{proposition}[Steady-state level in a room]
\label{prop:room}
\begin{equation}
  L_p(r) = L_W + 10\log_{10}\!\left(
    \frac{Q}{4\pi r^{2}} + \frac{4}{\mathcal{R}_{\mathrm{room}}}\right) .
  \label{eq:roomlevel}
\end{equation}
The two contributions are equal at the critical distance
\begin{equation}
  r_c = \sqrt{\frac{Q\,\mathcal{R}_{\mathrm{room}}}{16\pi}} .
  \label{eq:critdist}
\end{equation}
\end{proposition}

\begin{proof}[Sketch]
The direct term is \cref{eq:lwlp}. For the reverberant term, in the diffuse
field approximation the energy density is uniform and in steady state the power
absorbed equals the power injected: $W = \tfrac14 \bar\alpha S\, c_0\, E$ with
$E$ the energy density, whence $E = 4W/(c_0 S\bar\alpha)$; correcting for the
fact that the first reflection has already occurred replaces $S\bar\alpha$ by
$\mathcal{R}_{\mathrm{room}}$, and $p_{\mathrm{rms}}^2 = \rho c_0^2 E$ gives the $4/\mathcal{R}_{\mathrm{room}}$
term. Setting the two terms equal and solving for $r$ gives
\cref{eq:critdist}.
\end{proof}

\begin{finding}
\Cref{prop:room} has a consequence that no free-field figure can express:
beyond $r_c$, moving away from the machine buys nothing. The level is then set
by how absorbent the room is, not by how far away one stands. A flying display
does not make a quiet place near itself and a loud place far away; it raises
the level of the entire room to within a few decibels of the level at the user's
own ear.
\end{finding}

For an ordinary furnished room, $\bar\alpha\approx0.20$ and $r_c$ is under a
metre, which is smaller than the standoff at which such a machine is useful.
The user therefore sits close to the boundary between the two regimes and
receives both contributions (\cref{fig:room}).

\begin{figure}[tbp]
\centering
% Generated by paper/figures/gen_c.py. Do not edit by hand.
\def\cM{1.724}
\def\cMbat{0.288}
\def\cL{0.3827}
\def\cR{0.2460}
\def\cReach{0.629}
\def\cSpan{1.257}
\def\cKgap{1.100}
\def\cGap{0.0492}
\def\cSoverD{0.100}
\def\cKint{1.0252}
\def\cRcp{0.130}
\def\cCoMz{0.0275}
\def\cScreenOffz{0.1025}
\def\cScreenW{0.345}
\def\cScreenH{0.194}
\def\cArmr{8.0}
\def\cWall{0.66}
\def\cWallStress{0.083}
\def\cWallStiff{0.660}
\def\cWallMin{0.60}
\def\cDriver{stiffness}
\def\cSigmaWork{44}
\def\cSigmaAllow{350}
\def\cFbend{33}
\def\cFlimit{15.2}
\def\cZrotor{1.37}
\def\cZceil{1.03}
\def\cZoverR{5.57}
\def\cZcOverR{4.19}
\def\cKground{1.0020}
\def\cKceiling{1.0072}
\def\cKc{2.0}
\def\cCeilPct{0.7}
\def\cGroundPct{0.20}
\def\cKcSpanG{13}
\def\cCeiling{2.4}
\def\cBpf{34.5}
\def\cRpm{1036}
\def\cNb{2}
\def\cAbpf{-37.5}
\def\cLwa{61.8}
\def\cRroom{20.8}
\def\cRc{0.91}
\def\cStandoff{1.20}
\def\cSplEar{56.6}
\def\cSplOne{53.8}
\def\cSplFar{54.8}
\def\cTrotor{4.23}
\def\cVtip{26.7}
\def\cAnchor{82.9}
\def\cVi{3.01}
\def\cTfCross{21}
\def\cTf{30}
\def\cQmax{0.289}
\def\cMmot{61}
\def\cKtau{5.5}
\def\cEmot{0.30}
\def\cRoomL{5}
\def\cRoomW{4}
\def\cAlphaRoom{0.20}
%
% Eq:roomlevel for the design in the reference room: direct field, diffuse
% reverberant field, their sum, the critical distance and the user's ear.
\begin{tikzpicture}
  \begin{semilogxaxis}[paperplot, height=5.8cm,
      xlabel={distance from the machine $r$ [\si{\metre}]},
      ylabel={$L_p$ [\si{\decibel A}]},
      xmin=0.1, xmax=8, ymin=40, ymax=78,
      xtick={0.1,0.2,0.5,1,2,5}, xticklabels={0.1,0.2,0.5,1,2,5},
      legend pos=north east]
    \addplot[cB, dashed] table[x=r, y=direct] {results/fig/c_room.dat};
    \addlegendentry{direct, $L_W+10\log_{10}(Q/4\pi r^{2})$}
    \addplot[cE, dashed] table[x=r, y=reverb] {results/fig/c_room.dat};
    \addlegendentry{reverberant, $L_W+10\log_{10}(4/\mathcal{R}_{\mathrm{room}})$}
    \addplot[cA, line width=1.4pt] table[x=r, y=total] {results/fig/c_room.dat};
    \addlegendentry{total, \cref{eq:roomlevel}}
    \draw[construct] (axis cs:\cRc,40) -- (axis cs:\cRc,78);
    \node[lbl, cG, anchor=south east] at (axis cs:0.88,40.8) {$r_c=\SI{\cRc}{\metre}$};
    \addplot[only marks, mark=*, mark size=2.3pt, cA] coordinates {(\cStandoff,\cSplEar)};
    \node[lbl, cA, anchor=south west] at (axis cs:\cStandoff,\cSplEar+0.4)
      {user's ear, \SI{\cSplEar}{\decibel A}};
    \addplot[only marks, mark=o, mark size=2.3pt, cB] coordinates {(1,\cSplOne)};
    \draw[cB, line width=0.4pt] (axis cs:0.98,\cSplOne-0.3) -- (axis cs:0.82,51.2);
    \node[lbl, cB, anchor=north east] at (axis cs:0.84,51.4)
      {free field at \SI{1}{\metre}: \SI{\cSplOne}{\decibel A}};
    \node[lbl, cE, anchor=south, align=center] at (axis cs:4.2,54.9)
      {rest of the room};
  \end{semilogxaxis}
\end{tikzpicture}

\caption{\Cref{eq:roomlevel} for the design in the reference
$\cRoomL\times\cRoomW\times\SI{\cCeiling}{\metre}$ room with
$\bar\alpha=\cAlphaRoom$. The direct field falls by \SI{6}{\decibel} per
doubling of distance, the reverberant field does not fall at all, and the two
are equal at the critical distance $r_c=\SI{\cRc}{\metre}$ of
\cref{eq:critdist}. The user at \SI{\cStandoff}{\metre} is just past it and
hears \SI{\cSplEar}{\decibel A}, \SI{2.8}{\decibel} above the free-field figure
at one metre; the rest of the room settles at the reverberant level of about
\SI{\cSplFar}{\decibel A}.}
\label{fig:room}
\end{figure}

\subsection{Weighting and tonal content}

The A-weighting of IEC 61672 is
\begin{align}
  R_A(f) &= \frac{12194^2 f^4}
    {\big(f^2+20.6^2\big)\sqrt{\big(f^2+107.7^2\big)\big(f^2+737.9^2\big)}
     \big(f^2+12194^2\big)}, \nonumber\\
  A(f) &= 20\log_{10}R_A(f) + 2.00 .
  \label{eq:aweight}
\end{align}

The blade passage frequency is
\begin{equation}
  f_{\mathrm{BPF}} = \frac{N_b\,\Omega}{2\pi} = \frac{N_b\,\mathrm{rpm}}{60} .
  \label{eq:bpf}
\end{equation}

\begin{remark}[What the broadband number omits]
\Cref{eq:splmodel} predicts a broadband level. What makes a rotor
\emph{annoying} rather than merely loud is tonal content at
$f_{\mathrm{BPF}}$ and its harmonics. For the designs of \cref{sec:results},
$f_{\mathrm{BPF}}\approx\SI{\rBpf}{\hertz}$, which \cref{eq:aweight} discounts by
about \SI{38}{\decibel} (\cref{fig:aweight}). That is a genuine perceptual benefit of turning slowly
and it is also a warning: the acoustic energy has not gone anywhere. It has
moved to low frequency, where rooms absorb poorly, where Sabine's diffuse-field
assumption fails because the modal density is low, and where humans report a
throb rather than a hiss. Furthermore several rotors at slightly different
speeds beat at the difference frequency, which is perceptually worse than a
steady tone. None of this is in the A-weighted broadband figure, and it is the
most likely way for a built prototype to be unacceptable at a level the model
calls acceptable.
\end{remark}

The Schroeder frequency, above which a room behaves diffusely,
\begin{equation}
  f_S \approx 2000\sqrt{\frac{T_{60}}{V}} ,
  \label{eq:schroeder}
\end{equation}
is of order \SI{200}{\hertz} for a domestic room, which is precisely where the
first several blade harmonics lie. \Cref{prop:room} is therefore being applied
slightly outside its domain of validity for exactly the components that matter
most perceptually. This is stated as a limitation rather than repaired.

\begin{figure}[tbp]
\centering
% Generated by paper/figures/gen_c.py. Do not edit by hand.
\def\cM{1.724}
\def\cMbat{0.288}
\def\cL{0.3827}
\def\cR{0.2460}
\def\cReach{0.629}
\def\cSpan{1.257}
\def\cKgap{1.100}
\def\cGap{0.0492}
\def\cSoverD{0.100}
\def\cKint{1.0252}
\def\cRcp{0.130}
\def\cCoMz{0.0275}
\def\cScreenOffz{0.1025}
\def\cScreenW{0.345}
\def\cScreenH{0.194}
\def\cArmr{8.0}
\def\cWall{0.66}
\def\cWallStress{0.083}
\def\cWallStiff{0.660}
\def\cWallMin{0.60}
\def\cDriver{stiffness}
\def\cSigmaWork{44}
\def\cSigmaAllow{350}
\def\cFbend{33}
\def\cFlimit{15.2}
\def\cZrotor{1.37}
\def\cZceil{1.03}
\def\cZoverR{5.57}
\def\cZcOverR{4.19}
\def\cKground{1.0020}
\def\cKceiling{1.0072}
\def\cKc{2.0}
\def\cCeilPct{0.7}
\def\cGroundPct{0.20}
\def\cKcSpanG{13}
\def\cCeiling{2.4}
\def\cBpf{34.5}
\def\cRpm{1036}
\def\cNb{2}
\def\cAbpf{-37.5}
\def\cLwa{61.8}
\def\cRroom{20.8}
\def\cRc{0.91}
\def\cStandoff{1.20}
\def\cSplEar{56.6}
\def\cSplOne{53.8}
\def\cSplFar{54.8}
\def\cTrotor{4.23}
\def\cVtip{26.7}
\def\cAnchor{82.9}
\def\cVi{3.01}
\def\cTfCross{21}
\def\cTf{30}
\def\cQmax{0.289}
\def\cMmot{61}
\def\cKtau{5.5}
\def\cEmot{0.30}
\def\cRoomL{5}
\def\cRoomW{4}
\def\cAlphaRoom{0.20}
%
% The A-weighting of eq:aweight with the design's blade passage frequency
% and its harmonics, and the band below the Schroeder frequency where the
% diffuse-field room model does not hold.
\begin{tikzpicture}
  \begin{semilogxaxis}[paperplot, height=5.6cm,
      xlabel={frequency $f$ [\si{\hertz}]},
      ylabel={$A(f)$ [\si{\decibel}]},
      xmin=10, xmax=20000, ymin=-72, ymax=8,
      xtick={10,20,50,100,200,500,1000,2000,5000,10000,20000},
      xticklabels={10,,50,100,,500,1k,,5k,10k,},
      legend pos=south east]
    \fill[cD!8] (axis cs:10,-72) rectangle (axis cs:200,8);
    \node[lbl, cD, align=center, anchor=north] at (axis cs:45,6)
      {room not diffuse,\\$f<f_S\approx\SI{200}{\hertz}$};
    \addplot[cA, line width=1.2pt] table[x=f, y=A] {results/fig/c_aweight.dat};
    \addlegendentry{A-weighting, IEC 61672}
    \addplot[only marks, mark=*, mark size=2.2pt, cB] table[x=f, y=A] {results/fig/c_harm.dat};
    \addlegendentry{$k\,f_{\mathrm{BPF}}$, $k=1,\dots,6$}
    \node[lbl, cB, anchor=north west, align=left] at (axis cs:\cBpf,\cAbpf-1)
      {$f_{\mathrm{BPF}}=\SI{\cBpf}{\hertz}$\\$A=\SI{\cAbpf}{\decibel}$};
  \end{semilogxaxis}
\end{tikzpicture}

\caption{The A-weighting \cref{eq:aweight} with the design's blade passage
frequency and its first five harmonics. The fundamental at
\SI{\cBpf}{\hertz} is weighted down by \SI{37.5}{\decibel}, and every harmonic
up to the sixth lies in the shaded band below the Schroeder frequency
\cref{eq:schroeder}, where the diffuse-field model of \cref{prop:room} does not
hold. The A-weighted broadband figure is therefore least informative about
exactly the tonal components a listener notices.}
\label{fig:aweight}
\end{figure}
